\chapter{Math primer}

\section{What this primer is for}

The book has used a handful of mathematical ideas that come from
different parts of pure mathematics. None of them is essential to the
big picture --- a reader who skipped every formula would still
understand what the lens is doing --- but the precision they buy is
real, and a reader who wants to write down their own system in the
lens's vocabulary will want at least working command of them.

The four topics:

\begin{itemize}[leftmargin=1.5em,itemsep=2pt]
  \item \textbf{Order theory}: posets, lattices, and closure
    operators (used in Chapter~6).
  \item \textbf{Measure theory}: measurable spaces, measures, and
    Markov kernels (used throughout, especially Chapters~4 and~5).
  \item \textbf{Information theory}: entropy, mutual information, KL
    divergence, and effective information (used in Chapters~3, 7,
    and~8).
  \item \textbf{Category theory}: objects, morphisms, products, and
    monoidal categories (used lightly in Chapter~5 for
    $\boxtimes_\kappa$ and $\diamond_\rho$).
\end{itemize}

I will assume you have seen sets and functions, that you remember
what a probability density looks like, and that you can be persuaded
to read a definition slowly. I will not assume you have taken an
analysis class, a topology class, an algebra class, or a logic class.
Each section is self-contained.

The primer is meant to be readable end to end, but you can also use
it as a reference: skim until you hit the concept you need, read the
section it lives in, return to the chapter that prompted you.

\section{A short warm-up on sets and functions}

A \emph{set} is a collection of distinguishable objects. Examples:
$\{1, 2, 3\}$, the set of all real numbers $\mathbb{R}$, the set of
all $H \times W$ binary grids $\{0,1\}^{H\times W}$.

A \emph{function} $f : A \to B$ assigns to each element of $A$
exactly one element of $B$. We will use functions to describe
dynamics ($\Phi : X \to X$ maps a state to a next state), to describe
projections ($\pi : \prod_{p \in S} X_p \to Y$ maps a configuration
to a feature vector), and to describe maps between mathematical
structures (a closure operator is a function $2^P \to 2^P$).

The \emph{power set} of $A$, written $2^A$ or $\mathcal{P}(A)$, is
the set of all subsets of $A$. If $A = \{1, 2, 3\}$, then $2^A$ has
eight elements: the empty set, three singletons, three pairs, and
the full set. The power set will appear in order theory (the
canonical poset is $(2^A, \subseteq)$) and in measure theory (a
$\sigma$-algebra is a subset of the power set).

We will write $\mathbb{N}$ for the natural numbers $\{0, 1, 2,
\ldots\}$, $\mathbb{R}$ for the real numbers, and $\mathbb{R}_{\geq
0}$ for the non-negative reals.

That is the toolkit. The rest of this primer builds on it.

\section{Order theory}

Order theory is the study of relations that compare things. The book
uses it in Chapter~6 to talk about \emph{closure operators}, which
is how three of the four viability formalisms (autopoietic closure,
RAF, and arguably Markov-blanket) can be phrased.

\subsection*{Partially ordered sets (posets)}

A \emph{partial order} on a set $P$ is a binary relation $\sqsubseteq$
that is:

\begin{itemize}[leftmargin=1.5em,itemsep=2pt]
  \item \textbf{Reflexive}: $A \sqsubseteq A$ for every $A \in P$.
    Every element is comparable to itself.
  \item \textbf{Anti-symmetric}: if $A \sqsubseteq B$ and $B
    \sqsubseteq A$, then $A = B$. The order distinguishes distinct
    elements.
  \item \textbf{Transitive}: if $A \sqsubseteq B$ and $B \sqsubseteq
    C$, then $A \sqsubseteq C$.
\end{itemize}

A set $P$ with a partial order $\sqsubseteq$ is called a
\emph{partially ordered set} or \emph{poset}.

\paragraph{Canonical example: the power set.} For any set $X$, the
power set $2^X$ with the inclusion relation $\subseteq$ is a poset:
every set is a subset of itself; if $A \subseteq B$ and $B \subseteq
A$ then they are equal; subset-ness is transitive.

\paragraph{What ``partial'' means.} Some pairs of elements may not be
comparable. In $2^{\{1,2,3\}}$, the set $\{1\}$ is not a subset of
$\{2, 3\}$ and vice versa --- they are incomparable. A \emph{total}
order is one where every pair is comparable; the real numbers under
$\leq$ are totally ordered, but the power set of $\{1,2,3\}$ is not.

\paragraph{Other examples.} The natural numbers under divisibility:
$a \mid b$ iff $a$ divides $b$. Strings under prefix-of: $s
\sqsubseteq t$ iff $s$ is a prefix of $t$. Functions $\mathbb{R}
\to \mathbb{R}$ under pointwise order: $f \leq g$ iff $f(x) \leq
g(x)$ for all $x$.

\subsection*{Lattices, briefly}

Some posets have additional structure. A \emph{lattice} is a poset
in which every pair of elements has a least upper bound (their
\emph{join}, written $A \vee B$) and a greatest lower bound (their
\emph{meet}, written $A \wedge B$).

The power set $(2^X, \subseteq)$ is a lattice: $A \vee B = A \cup B$
(union is the least upper bound), $A \wedge B = A \cap B$
(intersection is the greatest lower bound). The natural numbers
under divisibility form a lattice: $a \vee b = \mathrm{lcm}(a, b)$,
$a \wedge b = \mathrm{gcd}(a, b)$.

The book uses lattices implicitly. Most of the time, ``the lattice
of subsets'' or ``the lattice of supporting sets'' just means the
power set with $\subseteq$. The fact that meets and joins exist is
what lets us talk about ``the smallest closed set containing $A$,''
which is the next concept.

\subsection*{Closure operators}

A \emph{closure operator} on a poset $(P, \sqsubseteq)$ is a function
$c : P \to P$ satisfying:

\begin{itemize}[leftmargin=1.5em,itemsep=2pt]
  \item \textbf{Extensivity}: $A \sqsubseteq c(A)$ for all $A$. The
    closure of an element is at least as big as the element.
  \item \textbf{Monotonicity}: if $A \sqsubseteq B$, then $c(A)
    \sqsubseteq c(B)$. Closure respects the order.
  \item \textbf{Idempotence}: $c(c(A)) = c(A)$. Closing twice is the
    same as closing once.
\end{itemize}

An element $A$ is \emph{closed} (under $c$) if $A = c(A)$.

\paragraph{Intuition.} A closure operator takes a candidate set and
``adds everything that follows.'' The result is the smallest closed
set that contains the input. If you put more in, you get more out
(monotonicity). If you close something already closed, nothing
changes (idempotence).

\paragraph{Examples that show up in the book.}

\begin{itemize}[leftmargin=1.5em,itemsep=2pt]
  \item \textbf{The convex hull.} On subsets of $\mathbb{R}^n$, the
    convex hull operator $c(A) = \mathrm{conv}(A)$ is a closure
    operator. The closed sets are exactly the convex sets.
  \item \textbf{Closure under a relation.} Given any binary relation
    $R$ on $X$, the operator ``add everything reachable by repeated
    application of $R$'' is a closure. The closed sets are the
    $R$-closed sets.
  \item \textbf{Autopoietic closure (Chapter~6).} On subsets of GoL
    cells, the operator ``add every cell whose next-state depends on
    the current set'' is a closure on the reaction-set lattice. The
    closed sets are the autopoietic units --- gliders, blocks, and
    related patterns.
  \item \textbf{RAF closure (Chapter~6).} On subsets of a chemical
    reaction network, the operator ``add every reaction whose
    reactants are present and whose catalyst is present'' is a
    closure. The closed sets are the RAFs.
\end{itemize}

\subsection*{Why closure operators matter for the book}

Conjecture~C1 in the framework says that each of the four viability
formalisms corresponds to a closure operator on some poset of
supporting sets. If true, viability would have a single abstract
form ($S$ is viable iff $S = c(S)$ for the formalism's $c$) and the
four formalisms would be choices of $c$.

The book takes the conjecture seriously without assuming it. The
language of closure operators is what makes the conjecture even
sayable.

\section{Measure theory}

Measure theory is the framework that lets us assign \emph{sizes} to
subsets of a space in a consistent way. The book uses it everywhere
the substrate state space appears: the substrate's $X$ is a
measurable space, the initial measure $\mu_0$ is a measure, the
population measure $\mu_t$ is a measure, and the variation operator
$V$ produces probability measures.

A reader's working model: ``measure'' is ``a generalisation of
volume.'' That model is enough for most of the book. The technical
content of measure theory is what makes the generalisation rigorous.

\subsection*{Sigma-algebras and measurable spaces}

The first surprise of measure theory is that you cannot always
assign sizes to \emph{every} subset of a space. Some subsets are too
weird --- the technical term is non-measurable --- and the only
sensible response is to refuse to give them a size at all.

A \emph{$\sigma$-algebra} on a set $X$ is a collection $\Sigma$ of
subsets of $X$ that:

\begin{itemize}[leftmargin=1.5em,itemsep=2pt]
  \item contains the empty set,
  \item is closed under complement (if $A \in \Sigma$ then $X
    \setminus A \in \Sigma$),
  \item is closed under countable union (if $A_1, A_2, \ldots \in
    \Sigma$, then $\bigcup_i A_i \in \Sigma$).
\end{itemize}

A set $A \in \Sigma$ is called \emph{measurable}. A pair $(X,
\Sigma)$ of a set and a $\sigma$-algebra is called a \emph{measurable
space}.

\paragraph{Don't panic.} For every space the book cares about, there
is a standard $\sigma$-algebra and you do not need to think about it
explicitly:

\begin{itemize}[leftmargin=1.5em,itemsep=2pt]
  \item For finite $X$ (a GoL grid of fixed size, an Avida program
    space of fixed instruction count), use the power set $\Sigma =
    2^X$. Every subset is measurable.
  \item For $\mathbb{R}^n$, use the Borel $\sigma$-algebra ---
    informally, the smallest $\sigma$-algebra that contains every
    open ball. Every subset you can describe constructively is in
    it.
  \item For Lenia's $[0,1]^L$ (a product of $|L|$ copies of $[0,1]$),
    use the product Borel $\sigma$-algebra.
\end{itemize}

The book's only use of the $\sigma$-algebra is to make ``measurable
function'' a precise notion: $f : X \to Y$ is measurable if the
preimage of every measurable subset of $Y$ is a measurable subset of
$X$. In practice, every function written down using continuous
operations is measurable. The notion is there to exclude
pathological cases.

\subsection*{Measures}

A \emph{measure} on a measurable space $(X, \Sigma)$ is a function
$\mu : \Sigma \to [0, \infty]$ such that:

\begin{itemize}[leftmargin=1.5em,itemsep=2pt]
  \item $\mu(\emptyset) = 0$ (the empty set has zero size).
  \item $\mu$ is \emph{countably additive}: for any disjoint
    sequence $A_1, A_2, \ldots$ of measurable sets, $\mu(\bigcup_i
    A_i) = \sum_i \mu(A_i)$.
\end{itemize}

The number $\mu(A)$ is the \emph{size} of $A$. Different measures
encode different notions of size.

\paragraph{Familiar measures.}

\begin{itemize}[leftmargin=1.5em,itemsep=2pt]
  \item On $\mathbb{R}^n$, \emph{Lebesgue measure} is the
    generalisation of length / area / volume. A unit ball in
    $\mathbb{R}^3$ has Lebesgue measure $\frac{4}{3}\pi$.
  \item \emph{Counting measure} on a finite or countable set:
    $\mu(A) = |A|$.
  \item \emph{Dirac measure} $\delta_x$ at a point $x \in X$:
    $\delta_x(A) = 1$ if $x \in A$, else $0$.
  \item \emph{Probability measures}: measures with $\mu(X) = 1$. The
    standard probabilistic distributions are all probability
    measures (a Gaussian on $\mathbb{R}$, the Bernoulli on
    $\{0,1\}$, etc.).
\end{itemize}

\paragraph{The two spaces of measures the book uses.}

\begin{itemize}[leftmargin=1.5em,itemsep=2pt]
  \item $\mathcal{P}(X)$: probability measures on $X$. The codomain
    of variation operators: $V : X \times \Theta \to \mathcal{P}(X)$.
    A probability measure represents a distribution over next
    states.
  \item $\mathcal{M}_+(X)$: finite positive measures on $X$. These
    are like probability measures but the total mass can be less
    than one. The codomain of viability filters: $F : \mathcal{M}_+(X)
    \to \mathcal{M}_+(X)$. The total mass can shrink (extinction is
    allowed) without leaving the space.
\end{itemize}

\paragraph{Why $\mathcal{P}(X)$ vs $\mathcal{M}_+(X)$ matters.} The
distinction is not bookkeeping. Variation produces normalised
distributions (the total probability of generating \emph{some}
variant is one). Viability can reduce the population. Conflating the
two would mean that viability could not extinguish a population,
which breaks the framework's commitment to allow extinction.

\subsection*{Pushforward}

If $\mu$ is a measure on $X$ and $f : X \to Y$ is a measurable
function, the \emph{pushforward} $f_* \mu$ is a measure on $Y$
defined by
\[
  (f_*\mu)(B) \;=\; \mu(f^{-1}(B)),
\]
i.e.\ the size of a subset $B \subseteq Y$ under $f_*\mu$ is the
size of its preimage in $X$ under $\mu$. The intuition: $f$ moves
mass from $X$ to $Y$; pushforward asks how much mass lands in each
subset of $Y$.

The book uses pushforward in two places. The lifted substrate
operator $\widetilde\Phi(x, \mu) = (\Phi(x), \Phi_* \mu)$ uses
pushforward to move the population measure along the substrate
dynamics. And the entity scale projection $\pi$ pushes forward the
substrate measure on a supporting set to a measure on the feature
space.

\subsection*{Markov kernels}

A \emph{Markov kernel} from $X$ to $Y$ is a measurable map
\[
  K : X \to \mathcal{P}(Y).
\]
Equivalently, $K$ takes each $x \in X$ to a probability
distribution $K(x) \in \mathcal{P}(Y)$. Sample from $K(x)$ to get a
value in $Y$.

Markov kernels are the standard mathematical object for representing
\emph{stochastic transitions}: from the current state $x$, you
sample the next state from the distribution $K(x)$.

\paragraph{Variation is a Markov kernel.} The book's variation
operator $V : X \times \Theta \to \mathcal{P}(X)$ is a Markov kernel
from $X \times \Theta$ to $X$. Deterministic variation is the
special case where $V(x, \theta)$ is a Dirac measure: all the
probability mass is on a single point.

\paragraph{Substrate dynamics is a Markov kernel.} A stochastic
substrate dynamics $\Phi : X \to \mathcal{P}(X)$ is a Markov kernel.
A deterministic dynamics $\Phi : X \to X$ is the special case where
the codomain is the set of Dirac measures.

\paragraph{Markov chains are iterated Markov kernels.} A
discrete-time Markov chain on a finite state space $X$ is determined
by a Markov kernel $P : X \to \mathcal{P}(X)$. We saw this in
Chapter~8 when computing effective information: the transition
matrix $P_{ij}$ is the value $P(i)(\{j\})$, i.e.\ the probability of
transitioning from $i$ to $j$.

\section{Information theory}

Information theory is the mathematical framework Claude Shannon
introduced in 1948 for measuring information. The book uses it in
Chapter~3 to set up causal emergence, in Chapter~7 to describe
information-theoretic observers, and in Chapter~8 to write down
effective information explicitly.

The fundamental object is \emph{entropy}, a number that captures
how spread out a probability distribution is.

\subsection*{Entropy}

Let $X$ be a random variable taking values in a finite set
$\mathcal{X}$ with probability $p(x) = \Pr[X = x]$. The
\emph{Shannon entropy} of $X$ is
\[
  \mathbf{H}(X) \;=\; -\sum_{x \in \mathcal{X}} p(x) \log_2 p(x),
\]
with the convention $0 \log 0 = 0$.

Properties:

\begin{itemize}[leftmargin=1.5em,itemsep=2pt]
  \item $\mathbf{H}(X) \geq 0$. Entropy is non-negative.
  \item $\mathbf{H}(X) = 0$ iff $X$ is deterministic (puts all mass
    on one point).
  \item $\mathbf{H}(X) \leq \log_2 |\mathcal{X}|$, with equality iff
    $X$ is uniform over $\mathcal{X}$.
  \item Entropy is measured in \emph{bits} when the logarithm is
    base 2, and in \emph{nats} when natural log is used. The book
    uses bits.
\end{itemize}

\paragraph{Intuition.} Entropy measures uncertainty. A fair coin
has $\mathbf{H} = 1$ bit: one bit of uncertainty. A two-headed coin
has $\mathbf{H} = 0$: no uncertainty. A fair die has $\mathbf{H} =
\log_2 6 \approx 2.58$ bits.

\paragraph{The $-p \log p$ shape.} The expression $-p \log p$ peaks
at $p = 1/e$ and tends to zero as $p \to 0$ or $p \to 1$. Summing
this over all values of a distribution gives entropy: high when many
values are about-equally-likely, low when one value dominates.

\subsection*{Joint and conditional entropy}

For two random variables $X, Y$ with joint distribution $p(x, y)$,
the \emph{joint entropy} is
\[
  \mathbf{H}(X, Y) \;=\; -\sum_{x, y} p(x, y) \log_2 p(x, y).
\]
The \emph{conditional entropy} of $Y$ given $X$ is
\[
  \mathbf{H}(Y \mid X) \;=\; \mathbf{H}(X, Y) - \mathbf{H}(X),
\]
the additional uncertainty in $Y$ once $X$ is known.

\subsection*{Mutual information}

The \emph{mutual information} between $X$ and $Y$ is
\[
  I(X; Y) \;=\; \mathbf{H}(X) + \mathbf{H}(Y) - \mathbf{H}(X, Y)
  \;=\; \mathbf{H}(Y) - \mathbf{H}(Y \mid X).
\]
Equivalently:
\[
  I(X; Y) \;=\; \sum_{x, y} p(x, y) \log_2 \frac{p(x, y)}{p(x)
  p(y)}.
\]
Properties:

\begin{itemize}[leftmargin=1.5em,itemsep=2pt]
  \item $I(X; Y) \geq 0$, with equality iff $X$ and $Y$ are
    independent.
  \item $I(X; Y) = I(Y; X)$ (symmetric).
  \item $I(X; Y) \leq \min\{\mathbf{H}(X), \mathbf{H}(Y)\}$, with
    equality iff one of $X, Y$ is a deterministic function of the
    other.
\end{itemize}

\paragraph{Intuition.} Mutual information is how much knowing one
variable tells you about the other. If $X$ and $Y$ are independent,
knowing $X$ tells you nothing about $Y$ and $I(X; Y) = 0$. If $Y =
f(X)$ deterministically, knowing $X$ tells you $Y$ exactly and $I(X;
Y) = \mathbf{H}(Y)$.

\paragraph{Where the book uses it.} The entity definition in
Chapter~4 uses conditional mutual information to express
informational closure: most of the predictive information about an
entity's future state should be inside its supporting set rather
than outside. The formula in the paper is
\[
  I(\pi(\mathbf{x}_S^{t+\tau});\,
  \pi(\mathbf{x}_S^{t}) \mid \mathbf{x}_{\bar S}^{t})
  \;\geq\; I(\pi(\mathbf{x}_S^{t+\tau});\, \mathbf{x}_S^{t})
  - \varepsilon,
\]
which reads: the information that the supporting set's past carries
about its future (conditional on the outside) is large compared to
the information the inside carries unconditionally. In plain
language: the entity is informationally insulated from its
surroundings.

\subsection*{KL divergence}

The \emph{Kullback--Leibler divergence} between two distributions
$p$ and $q$ on the same space is
\[
  D_{\mathrm{KL}}(p \,\|\, q) \;=\; \sum_x p(x) \log_2 \frac{p(x)}{
  q(x)},
\]
again with $0 \log 0 = 0$ and a divergence of $\infty$ if $p$ puts
mass anywhere $q$ does not.

Properties:

\begin{itemize}[leftmargin=1.5em,itemsep=2pt]
  \item $D_{\mathrm{KL}}(p \,\|\, q) \geq 0$ (Gibbs's inequality),
    with equality iff $p = q$.
  \item $D_{\mathrm{KL}}$ is \emph{not} symmetric in general:
    $D_{\mathrm{KL}}(p \,\|\, q) \neq D_{\mathrm{KL}}(q \,\|\, p)$.
    It is therefore not a distance in the metric sense.
  \item $I(X; Y) = D_{\mathrm{KL}}(p(x,y) \,\|\, p(x) p(y))$. Mutual
    information is the KL divergence between the joint distribution
    and the product of marginals.
\end{itemize}

\paragraph{Intuition.} $D_{\mathrm{KL}}(p \,\|\, q)$ is how much
information you lose if you use the wrong distribution $q$ to encode
data from $p$. Smaller is better; zero means $q$ matches $p$.

\subsection*{Effective information}

We can now write the effective information formula from Chapter~8
explicitly. Let $M = (X, P)$ be a discrete-time Markov chain on a
finite state space $X$ of size $n$, with transition matrix $P$.
Define the \emph{marginal effect distribution} $\bar P_\cdot =
\frac{1}{n} \sum_{i=1}^n P_{i, \cdot}$ --- the row average of $P$.

The \emph{effective information} of $M$ under uniform intervention
is
\[
  \mathrm{EI}(M) \;=\; \frac{1}{n} \sum_{i=1}^n D_{\mathrm{KL}}(P_{
  i, \cdot} \,\|\, \bar P_\cdot).
\]
Equivalently,
\[
  \mathrm{EI}(M) \;=\; \mathbf{H}(\bar P_\cdot) - \frac{1}{n}
  \sum_i \mathbf{H}(P_{i, \cdot}).
\]

\paragraph{Reading the formula.} The first form says: take each
row's deviation from the average row (in KL divergence), and
average. The second form says: take the entropy of the average row,
minus the average row entropy. Both are equal.

\paragraph{Two extreme cases.}

\begin{itemize}[leftmargin=1.5em,itemsep=2pt]
  \item All rows are identical. Then $\bar P_\cdot = P_{i, \cdot}$
    for every $i$, so $D_{\mathrm{KL}}(P_{i, \cdot} \,\|\, \bar
    P_\cdot) = 0$ and $\mathrm{EI} = 0$. The chain has no causal
    power: every starting state goes to the same distribution of
    next states.
  \item Each row is a Dirac on a different next state, and all next
    states are distinct. Then each row has entropy $0$ (deterministic)
    and the average row is uniform over $n$ values (entropy $\log_2
    n$). So $\mathrm{EI} = \log_2 n - 0 = \log_2 n$, the maximum.
\end{itemize}

A system has high effective information when interventions matter
--- when different starting states reliably lead to different and
distinguishable outcomes.

\section{Category theory}

Category theory is the most abstract of the four topics, and the
book uses it most lightly. It enters in Chapter~5 to justify the
substrate-composition operator $\boxtimes_\kappa$ and the
variation-channel composition operator $\diamond_\rho$ as
``symmetric monoidal'' constructions. You can read the whole book
without learning what those words mean. This section gives the
working vocabulary so that you can recognise the references when
you see them.

\subsection*{Categories}

A \emph{category} $\mathcal{C}$ consists of:

\begin{itemize}[leftmargin=1.5em,itemsep=2pt]
  \item A collection of \emph{objects} (think: sets, spaces,
    substrates).
  \item For each pair of objects $A, B$, a collection of
    \emph{morphisms} $A \to B$ (think: functions, maps,
    transformations).
  \item For each object $A$, an \emph{identity morphism} $\mathrm{id
    }_A : A \to A$.
  \item A \emph{composition operation} $\circ$: given $f : A \to B$
    and $g : B \to C$, you get $g \circ f : A \to C$.
\end{itemize}

Subject to two axioms: composition is associative ($h \circ (g
\circ f) = (h \circ g) \circ f$) and identities behave as expected
($f \circ \mathrm{id}_A = f = \mathrm{id}_B \circ f$).

\paragraph{Canonical examples.}

\begin{itemize}[leftmargin=1.5em,itemsep=2pt]
  \item $\mathbf{Set}$: objects are sets, morphisms are functions,
    composition is the usual composition of functions.
  \item $\mathbf{Meas}$: objects are measurable spaces, morphisms
    are measurable functions.
  \item $\mathbf{Top}$: objects are topological spaces, morphisms
    are continuous functions.
\end{itemize}

A category is a vocabulary for talking about \emph{things plus
relationships between things}, abstracting away from what the things
are made of. The same theorems can apply to sets, spaces, groups,
substrates --- whatever fits the schema.

\subsection*{Products}

A \emph{product} of two objects $A$ and $B$ in a category is an
object $A \times B$ together with two projection morphisms
\[
  \pi_A : A \times B \to A, \quad \pi_B : A \times B \to B,
\]
such that for any other object $C$ with morphisms $f : C \to A$ and
$g : C \to B$, there is a unique morphism $\langle f, g\rangle : C
\to A \times B$ making the obvious diagram commute.

\paragraph{Why this is interesting.} In $\mathbf{Set}$, products
are just Cartesian products. In $\mathbf{Top}$, products carry the
product topology. In $\mathbf{Meas}$, products carry the product
$\sigma$-algebra. The categorical definition unifies them: ``the
right product'' is whatever satisfies this universal property in
the category at hand.

The substrate composition $\mathbf{X}_1 \boxtimes \mathbf{X}_2$ in
Chapter~5 is the categorical product in a category of substrates
with substrate-preserving morphisms.

\subsection*{Monoidal categories}

A \emph{monoidal category} is a category equipped with a binary
operation $\otimes$ on objects (and morphisms) that behaves like
multiplication: there is a unit object, the operation is
associative up to isomorphism, and these properties satisfy
coherence axioms.

A monoidal category is \emph{symmetric} if there is a natural
isomorphism $A \otimes B \cong B \otimes A$.

\paragraph{Why this matters for the book.} The substrate
composition operator $\boxtimes$ in its free form (no coupling) is
a symmetric monoidal product on the category of substrates. The
coupled form $\boxtimes_\kappa$ is a symmetric monoidal product on
the category of substrates equipped with coupling data. The
variation channel composition $\diamond_\rho$ has the same
structure on the category of variation channels.

The payoff: associativity up to canonical isomorphism. You can write
$(\mathbf{X}_1 \boxtimes_{\kappa_{12}} \mathbf{X}_2)
\boxtimes_{\kappa_{23}} \mathbf{X}_3$ or $\mathbf{X}_1
\boxtimes_{\kappa_{12}} (\mathbf{X}_2 \boxtimes_{\kappa_{23}}
\mathbf{X}_3)$ and the two are canonically the same composite. You
do not need to choose a parenthesisation when building
biological-atmospheric-cultural composites.

\subsection*{Operads and wiring diagrams}

A more advanced construction: \emph{operads of wiring diagrams} are
categorical structures that formalise ``boxes with wires plugged
into each other.'' They are the natural home for systems built by
coupling components, and they are the framework's formal foundation
for the substrate and variation composition operators.

The book does not require you to understand operads. The reference
in Chapter~5 is for completeness: if you know the operad-of-wiring-
diagrams literature, the framework's compositions are operad maps
in that sense. If you do not, no harm done.

\subsection*{Functors and natural transformations, briefly}

Two more pieces of categorical vocabulary that occasionally appear
in the artificial-life literature:

\begin{itemize}[leftmargin=1.5em,itemsep=2pt]
  \item A \emph{functor} $F : \mathcal{C} \to \mathcal{D}$ is a map
    between categories that takes objects to objects, morphisms to
    morphisms, and preserves composition and identities. Functors
    are ``structure-respecting maps between categories.''
  \item A \emph{natural transformation} $\eta : F \Rightarrow G$
    between two functors is a family of morphisms $\eta_A : F(A)
    \to G(A)$, one for each object, that respect morphisms in the
    obvious sense.
\end{itemize}

These rarely show up explicitly in the book. They are how
categorical formulations express ``a uniform way of relating one
construction to another.''

\section{Putting it together}

We have surveyed four mathematical regions. They are connected to
the book like this:

\begin{itemize}[leftmargin=1.5em,itemsep=2pt]
  \item \textbf{Order theory} $\Rightarrow$ closure operators
    $\Rightarrow$ viability formalisms in Chapter~6 $\Rightarrow$
    Conjecture~C1.
  \item \textbf{Measure theory} $\Rightarrow$ measurable spaces and
    measures $\Rightarrow$ substrate state spaces, population
    measures, Markov kernels $\Rightarrow$ the lens's signatures
    throughout.
  \item \textbf{Information theory} $\Rightarrow$ entropy, mutual
    information, KL divergence $\Rightarrow$ entity informational
    closure (Chapter~4), effective information and causal emergence
    (Chapters~3, 8).
  \item \textbf{Category theory} $\Rightarrow$ monoidal categories
    $\Rightarrow$ $\boxtimes_\kappa$ and $\diamond_\rho$
    compositions (Chapter~5).
\end{itemize}

You do not need to internalise everything in this primer to use the
book. A reader who can write a probability density and read a
function signature can read the lens fluently after this primer
with occasional reference back to specific sections.

What the primer is for is the moment when a chapter writes down
something like $V : X \times \Theta \to \mathcal{P}(X)$ or
$\mathrm{EI}(M^\Pi) - \mathrm{EI}(M)$ or ``the closure operator on
the lattice of supporting sets,'' and you want to know what those
symbols are doing. Flip back here. The full machinery is light, but
it is precise, and the precision is what makes the lens work.

\section*{Further reading}

For information theory: Cover \& Thomas, \textit{Elements of
Information Theory}, is the standard textbook and is accessible to
a reader with calculus.

For measure theory: Pollard's \textit{User's Guide to Measure
Theoretic Probability} is short and conceptual. Folland's
\textit{Real Analysis} is the standard reference and is harder.

For order theory and lattices: Davey \& Priestley,
\textit{Introduction to Lattices and Order}, is the standard.
Closure operators are covered in the first few chapters.

For category theory: Spivak's \textit{Category Theory for the
Sciences} is accessible for applied readers. Awodey's
\textit{Category Theory} is more advanced.

For Hoel-style causal emergence: the source papers by Hoel,
Tononi, and Albantakis, and follow-up work by Klein and Hoel and
by Comolatti and Hoel, contain the formalism and the principal
applications.
